\chapter{Theoretical Background} \label{ch:chapter_04}
\section{Thermal Control in Engineering}
In engineering applications, thermal control is decisive to optimize the 
performance of the system and its components, while keeping them within temperature limits. In some systems 
and elements, such as engines, electric components or aeroespace systems, the implementation of a \acrfull{tcs} 
is critical because thermal variations can produce a negative 
impact on their performance and structural safety \cite{th_handbook}.

Methods such as convection, conduction and radiation can be implemented on Earth's surface to study heat 
transfer phenomena. However, the lack of an atmosphere in space, limits heat transfer mechanisms to conduction 
in solids and thermal radiation into space \cite{spacethctrl}. 

The primary objective of a Spacecraft \acrshort{tcs} is to ensure that all onboard components operate within 
their specified temperature ranges during all mission 
phases, while minimizing resource consumption (mass and power)\cite{tcs_slides0}. 
This is achieved through a careful balance of internal heat dissipation and external environmental fluxes, utilizing 
both passive and active thermal control technologies \cite{th_handbook}.


Those aspects are critical because, unlike terrestrial applications, the space environment is characterized by extreme conditions, 
such as vacuum, a background temperature of 3K, and rapid, intense thermal cycling due to 
exposure to sunlight and Earth's shadow \cite{tcs_slides0,spacethctrl}.


Failure to comply with these thermal limits can lead to equipment malfunctions, accelerated material 
degradation, reduced reliability, structural damage from thermal cycling, or even catastrophic failure.

The objective of thermal control is to prevent those situations, ensuring operational reliability and 
enhancing the efficiency of the system \cite{hmt_f&a}. Some key requirements for 
the \acrshort{tcs} include \cite{tcs_slides1}:
\begin{itemize}[label={\scriptsize\raisebox{0.5ex}{\textbullet}}]
	\item Temperature Ranges: maintaining specific operational and survival temperature ranges 
	for components: electronics, batteries or infrared detectors
	\item Thermal Gradients and Stability: controlling temperature gradients and stability is essential 
	for high-precision systems such as cameras, optical benches, and CCD sensors to 
	avoid structural distortions and measurement errors
\end{itemize}
 
\section{Space Thermal Environment}

Spacecraft thermal control is a process of energy management in which environmental heating plays a major role 
and therefore, has a major impact in the spacecraft design. There is no convection in space due to the absence of an atmosphere,
being radiation the only possible method to dissipate heat in space. \newline
Nevertheless, heat conduction between the structure and the components of the spacecraft still occurs. \cite{th_handbook}

\section{Heat Transfer Mechanisms}

Heat can be defined as the form of energy that can be transferred from one system to another 
as a result of temperature difference. The science that deals with the determination of the rates of such 
energy transfers is heat transfer \cite{hmt_f&a}.
Conduction, convection and radiation are the fundamental mechanisms of heat transfer. However, as it is mentioned earlier,
there is no convection in space systems, leaving conduction and radiation as the predominant mechanisms. \cite{th_handbook}

% EJEMPLO DE COMO PONER UNA IMAGEN
%\begin{figure}[H]
%	\centering
%	\includegraphics[width=100mm]{imag_07.png}
%	\caption{ Spectral distribution of the emissive power of a blackbody at different temperatures.
%The peak shifts to shorter wavelengths as temperature increases, according to Planck’s Law, \cite{heattransfer}} % Pie de imagen
%	\label{fig:001_01}
%\end{figure}

\subsection{Conduction Heat Transfer}

Heat transfer conduction is the mechanism by which thermal energy is transmitted 
through solids or between solids in contact due to the presence of a temperature gradient, without any 
macroscopic motion of the substance. This process is governed by microscopic interactions between 
particles and, in metallic materials, the transport of free electrons \cite{hmt_f&a}.

From a macroscopic perspective, heat conduction is described by Fourier’s law:
\begin{equation}
q = -k A \frac{dT}{dx}
\label{eq:Fourier}
\end{equation}

where: 
\begin{itemize}[label={\scriptsize\raisebox{0.5ex}{\textbullet}}]
	\item $q$ is the heat transfer rate (W) 
	\item $k$ is the thermal conductivity (W/mK) 
	\item $A$ is the cross-sectional area (m$^2$) 
	\item $\frac{dT}{dx}$ is the temperature gradient (K/m). 
\end{itemize}

The negative sign in the temperature gradient indicates that heat flows from regions of higher 
temperature to lower temperature \cite{hmt_f&a, spacethctrl}.

Under steady-state conditions and simple geometries, 
Fourier’s law can be integrated to obtain analytical expressions for heat transfer. 
For instance, in a plane wall with constant thermal conductivity, the heat flux is directly proportional to the temperature
difference between the boundaries and inversely proportional to the wall thickness \cite{hmt_f&a}. \newline
\autoref{fig:001_07} illustrates one-dimensional heat conduction through a flat wall, consistent with the
formulation of Fourier's law in \autoref{eq:Fourier}. 

\begin{figure}[H]
	\centering
	\includegraphics[width=100mm]{imag_07.png}
	\caption{One-dimensional heat conduction accross the wall, with thickness $\Delta x$ ,whose surfaces are at 
temperatures  T1  and  T2  (T1  >  T2), respectively \cite{spacethctrl}} % Pie de imagen
	\label{fig:001_07}
\end{figure}

For more general cases involving spatial and temporal variations, the heat conduction process is 
governed by the heat diffusion equation, expressed in cartesian coordinates as:

\begin{equation}
\rho c_p \frac{\partial T}{\partial t} = \nabla \cdot (k \nabla T) + \dot{q}
\end{equation}

where $\rho$ is the density (kg/m$^3$), $c_p$ is the specific heat (J/kgK), 
and $\dot{q}$ represents internal heat generation per unit volume (W/m$^3$). 
This equation forms the basis for thermal modelling in both analytical and numerical approaches, 
including finite difference and nodal methods \cite{hmt_f&a, thermal_eng_man, ecss_man}.

An additional useful concept in conduction analysis is thermal resistance, which simplifies 
the study of heat transfer in complex systems. 
For a homogeneous layer, the conductive thermal resistance is defined as:

\begin{equation}
R_{cond} = \frac{L}{kA}
\end{equation}

where $L$ is the thickness of the material. \newline
In multilayer structures, individual resistances can be combined in series or parallel, allowing the overall 
heat transfer to be evaluated using an approach analogous to electrical circuits \cite{hmt_f&a}.

In the context of spacecraft thermal control, conduction plays a fundamental role in the 
internal redistribution of heat. Conductive paths are essential to transport heat generated 
by onboard components towards external radiative surfaces. Therefore, the design of conductive 
interfaces, material selection, and contact conductance are critical factors in ensuring proper 
thermal management \cite{spacethctrl, th_handbook, ecss_man}.

Furthermore, in thermal modelling tools such as ESATAN, conduction is 
typically represented through thermal conductances connecting discrete nodes within a network.
These conductances depend on material properties and geometry and are used to solve the energy balance 
equations governing the thermal behaviour of the system \cite{thermal_eng_man}.

\subsection{Radiation Heat Transfer}
%\begin{figure}[H]
%	\centering
%	\includegraphics[width=100mm]{imag_07.png}
%	\caption{ Spectral distribution of the emissive power of a blackbody at different temperatures.
%The peak shifts to shorter wavelengths as temperature increases, according to Planck’s Law, \cite{heattransfer}} % Pie de imagen
%	\label{fig:001_01}
%\end{figure}
\section{Mathematical Modeling of Thermal Systems}

Mathematical modeling is essential to predict the thermal behavior of space systems and design
effective thermal control strategies. 

\subsection{Disadvantages of the Finite Element Method (FEM)}



\subsection{Lumped Parameter Method (LPM)}

The \acrfull{lpm} is a discretization technique that simplifies a continuous physical system 
into a finite number of discrete elements called nodes. 
The fundamental assumption is that each node is isothermal, meaning its temperature is uniform throughout its entire volume. 
This allows the system to be solved as a network of interconnected thermal resistances and 
capacitances, analogous to an electrical circuit, 
where temperature equates to voltage and heat flow to current flow. \cite{thermal_eng_man}

\autoref{fig:001_04} shows a simplified thermal system and its equivalent nodal network representation:

\begin{figure}[H]
	\centering
	\includegraphics[width=130mm]{imag_04.png}
	\caption{Simplified heat transfer system and its corresponding thermal network \cite{lpm_elec}} % Pie de imagen
	\label{fig:001_04}
\end{figure}

The behaviour of the network can be described as a system of differential equations whose coefficients are grouped values of the nodes. The mathematical expression for the nodal energy balance equation is:

\begin{equation}
    C_i \frac{dT_i}{dt} = Q_i - \sum_{j=1}^{n} GL_{ij}(T_i - T_j) - \sum_{j=1}^{n} GR_{ij}(T_i^4 - T_j^4)
\end{equation}
where:

\begin{itemize}[label={\scriptsize\raisebox{0.5ex}{\textbullet}}]
	\item $C_i$ [J/K]: thermal capacitance of node $i$ ($m_i\ c_p,_i$); stores/releases heat as $T_i$ changes.  %El dolar se usa para ponerlo en formato formula

	\item $T_i$ [K]: temperature of node $i$.

	\item $t$ [s]: time.

	\item $GL_{ij}$ [W/K]: conductive conductance between nodes $i$ and $j$. The value is zero if there is not a conductive path.

	\item $GR_{ij}$ [$W/K^4$]: radiative coupling factor between $i$ and $j$ in the non-linear $T^4$ form; it embeds
$\sigma$, view factors and emissivities for diffuse–gray surfaces.

	\item $Q_i$ [W]: net heat source applied to node $i$ (internal dissipation, heaters, and absorbed external fluxes as modelled in the network).

	\item $\sum_{j=1}^n$: sum over all nodes $j$, which are thermally connected to $i$. The terms $GL_{ij}(T_i-T_j)$ and $GR_{ij}(T_i^4-T_j^4)$ contribute as inflow when the expression in parentheses is negative.
\end{itemize}


The discrete approximation simplifies the numerical solution of the thermal problem by reducing 
a continuous system, governed by differential equations, to a finite set of equations 
defined over a limited number of nodes. In steady-state conditions, this system becomes algebraic 
and can be solved using standard numerical techniques, whereas in transient analysis it remains a 
system of differential equations \cite{thermal_eng_man}.

Once the energy balance is defined for each node, the complete model is obtained by assembling 
all the equations into a coupled system. This system describes the temporal evolution of the nodal 
temperatures, taking into account both internal heat sources and the thermal interactions between nodes. \cite{thermal_eng_man}

In this framework, thermal systems are represented as networks of nodes connected through thermal 
conductances. These conductances account for the main heat transfer mechanisms, such as conduction 
through structural elements and radiation exchange between surfaces. Each node is characterised by 
its thermal capacitance and exchanges heat with the rest of the system, enabling the prediction of 
both steady-state and transient behaviour \cite{thermal_eng_man, spacethctrl}.

The accuracy of this method depends on the discretisation of the system. A reduced number of nodes 
may not capture the relevant temperature gradients, while an excessive number increases the 
computational cost without providing significant accuracy in the results. Therefore, the definition of 
the nodes represents a compromise between accuracy and efficiency. \cite{ecss_man}

Therefore, due to its balance between accuracy and computational cost, this method is widely used in spacecraft
thermal modelling and forms the basis of tools such as ESATAN-TMS. \cite{ecss_man}
%Therefore, this method has a wide range of applications in the aeroespace industry, such as the thermal modeling of satellites, 
%where a balance between accuracy and efficiency is required.

\section{Software and Advanced Methods for Thermal Simulation}

\section{Analytical Capabilities of ESATAN-TMS}



